% Draft manuscript. Numbers come from runs/paper (paper_runs/) as of 4 Nov 2026: one seed per model.
% Segmentation and end-to-end results are the Colab output; late fusion (kfold_cv, split_gap, negative_control,
% learning_curves) was rerun on CPU. Figures are read from ../paper_runs/paper_figures/ (not in git); see README.md.
% If seeds 43-44 are run later, replace the one-seed results with mean +- SD (scripts/aggregate_results.py).
\documentclass[11pt]{article}

\usepackage[T1]{fontenc}
\usepackage[utf8]{inputenc}
\usepackage[a4paper,margin=2.5cm]{geometry}
\usepackage{amsmath}
\usepackage{booktabs}
\usepackage{graphicx}
\usepackage{subcaption}
\usepackage{xcolor}
\usepackage[numbers,sort&compress]{natbib}
\usepackage[hidelinks]{hyperref}

\graphicspath{{../paper_runs/paper_figures/}}
\newcommand{\todo}[1]{\textcolor{red}{[TODO: #1]}}
\newcommand{\ci}[2]{\,[#1, #2]}

\title{What does T2-weighted MRI add to glioma grading and WHO~2021 subtyping? A leakage-controlled benchmark of
segmentation, deep features and radiomics}
\author{Konstantinos Koutsikos\thanks{School of Electrical and Computer Engineering, National Technical
University of Athens, Greece. \todo{co-authors, e-mail}}}
\date{}

\begin{document}
\maketitle

\begin{abstract}
\noindent\textbf{Purpose.} Deep-learning pipelines that grade gliomas and predict their WHO~2021 subtype from MRI
are usually developed on small cohorts. There, label leakage and single-split evaluation can inflate results, and it
is often unclear what the images add beyond demographics. We built a leakage-controlled multimodal pipeline and
measured what each input and design choice contributes.

\noindent\textbf{Methods.} On T2-weighted MRI from the UCSF-PDGM dataset we benchmarked seven segmentation
configurations (SwinUNETR with and without self-supervised pretraining, SegResNet, two 3D U-Net configurations,
SegMamba at two learning rates) on 75 held-out patients. On a grade-balanced 162-patient cohort we trained
end-to-end multi-task classifiers and late-fusion models for three-way WHO grade, four-way WHO~2021 diagnosis, LGG
versus HGG and IDH status. Inputs were limited to the images, pre-operative clinical variables and radiomics.
Molecular and surgical variables, which define or follow the labels, were used only in a labelled reference
experiment. One nested patient-level split was fixed before training and every feature selector was fit on training
data only. We report bootstrap confidence intervals, pre-specified paired tests, repeated 5-fold cross-validation
(CV) and an analysis with ground-truth tumour masks. Post hoc, we tested validation--test gaps against chance
re-assignment of patients and ran a shuffled-label control.

\noindent\textbf{Results.} SegResNet reached the highest Dice (0.865, 95\% CI 0.828--0.899) and the lowest HD95
(11.7~mm). The best SegMamba and SwinUNETR runs matched its Dice but had significantly larger HD95, and
self-supervised pretraining gave no gain (Dice 0.842 vs 0.843). Molecular and surgical variables (IDH, 1p/19q,
MGMT, extent of resection) alone gave a CV diagnosis AUC of 0.93 without any imaging. Without them, clinical variables
plus radiomics reached a CV diagnosis AUC of $0.830 \pm 0.051$ against $0.652 \pm 0.045$ for clinical variables alone
with XGBoost, and 0.830 against 0.727 with logistic regression (both corrected $p < 0.01$). Deep image features gave a
similar gain on their own but none once radiomics were in the model (XGBoost $0.831 \pm 0.046$, $p = 0.99$), and
radiomics alone were not significantly better than deep features alone. Inside age bands, radiomics kept an LGG versus
HGG AUC of 0.66--0.81 against 0.45--0.62 for age and sex (logistic regression). Age and sex alone predicted IDH status
with a CV AUC of 0.85--0.90; imaging added nothing significant after correction and did not clearly beat age inside age
bands. On shuffled labels the pipeline scored 0.49--0.51, against up to 0.58 when features were selected before the
folds. For diagnosis, chance re-assignment of patients alone moved the gap between the two 25-patient splits by about
$\pm 0.2$ AUC (95\% range). An end-to-end image + clinical + radiomics model reached a test diagnosis AUC of 0.916
(0.837--0.985) but 0.619 on validation, and no pre-specified difference between end-to-end models was significant
after correction.

\noindent\textbf{Conclusions.} Without leakage, T2-weighted imaging adds to clinical variables for diagnosis through
either radiomics or deep features, which are largely redundant with each other. For grade, diagnosis and LGG versus HGG
it still beats age and sex inside age bands, so it is not merely a proxy for age; for IDH status we could not show
this. Results on a single 25-patient split are too variable to rank models. Code and the experiment protocol are
public.
\end{abstract}

\noindent\textbf{Keywords:} glioma, WHO~2021, radiogenomics, data leakage, segmentation benchmark, SegMamba,
self-supervised learning, radiomics, small data.

% ---------------------------------------------------------------------------------------------------------
\section{Introduction}

The 2021 WHO classification of central nervous system tumours made molecular markers part of the diagnosis of
diffuse gliomas: IDH status separates IDH-mutant astrocytomas and oligodendrogliomas from IDH-wildtype
glioblastomas, and 1p/19q codeletion defines oligodendroglioma~\cite{louis2021who}. Predicting these markers and the
WHO grade from pre-operative MRI (a ``virtual biopsy'') has been studied extensively with deep networks and
radiomics~\cite{chang2018ajnr,chang2018ccr,choi2021idh,patel2017mismatch}. Most such studies train on a few hundred
patients or fewer. In that regime, how a model is evaluated can move results as much as which model it is.
Leakage from inputs that encode the label, feature selection done before the split, overlapping training and test
patients across pipeline stages, and conclusions drawn from one small test set are all common in medical machine
learning~\cite{kapoor2023leakage,varoquaux2022failures,varoquaux2018cv}.

Typical pipelines chain three steps: they segment the tumour, extract features, and classify. Each step invites a
design choice whose value is rarely isolated:
\begin{itemize}
  \item state-space and transformer segmenters~\cite{xing2024segmamba,hatamizadeh2022swinunetr} against convolutional
        ones~\cite{myronenko2019segresnet,isensee2021nnunet}, with or without self-supervised
        pretraining~\cite{tang2022ssl};
  \item deep features from the segmentation network against hand-crafted radiomics~\cite{vangriethuysen2017pyradiomics};
  \item end-to-end multi-task fusion against gradient-boosted late fusion~\cite{chen2016xgboost}.
\end{itemize}
Clinical variables are a further source of ambiguity. Some, such as age, are known before surgery. Others, such as
IDH, 1p/19q and extent of resection, come from tissue or surgery, and IDH and 1p/19q define the diagnosis being
predicted.

We built one pipeline in which each of these choices can be switched independently, held every evaluation detail
fixed, and asked four questions on T2-weighted MRI:
\begin{enumerate}
  \item Which architecture segments the whole tumour best, and does self-supervised pretraining help?
  \item What does imaging add for grade, diagnosis, LGG versus HGG and IDH status beyond age and sex, and do deep
        features add to hand-crafted radiomics?
  \item How much would molecular and surgical inputs inflate the results?
  \item How reliable is a ranking of models on a 25-patient test split, and what do segmentation errors cost
        downstream?
\end{enumerate}
Our contributions are:
\begin{itemize}
  \item a segmentation benchmark under one matched budget, with paired statistics;
  \item an input ladder evaluated both end to end and by late fusion with repeated CV and pre-specified tests;
  \item a labelled leakage reference that quantifies how molecular inputs produce near-perfect numbers;
  \item a permutation test of validation--test gaps, and an analysis of how segmentation errors propagate to
        classification;
  \item a shuffled-label control of the whole late-fusion pipeline, which also measures what feature selection
        before the folds would have cost;
  \item an open, resumable implementation.
\end{itemize}

% ---------------------------------------------------------------------------------------------------------
\section{Materials and methods}

\subsection{Data}
We used the UCSF Preoperative Diffuse Glioma MRI dataset (UCSF-PDGM)~\cite{calabrese2022ucsf}: 495 adults with
histopathologically confirmed WHO grade 2--4 diffuse glioma, 3~T pre-operative MRI, expert tumour segmentations and
clinical, molecular and survival metadata. The dataset's 6 follow-up scans of these patients were excluded. We used
only the T2-weighted sequence, which is part of every glioma protocol, to study what a single structural sequence
carries. The segmentation target was the whole tumour, with all labelled compartments merged.

The cohort is dominated by grade 4 (396 of 495 patients). Classification therefore used a grade-balanced
162-patient sub-cohort: all 56 grade-2 and 43 grade-3 patients and 63 randomly drawn grade-4 patients, so 333 grade-4
patients were excluded. Its median age was 48 years (range 19--81), 38.9\% were female, and 81 tumours were IDH-mutant
(Table~\ref{tab:cohort}). Targets were:
\begin{itemize}
  \item three-way WHO grade;
  \item four-way WHO~2021 diagnosis: astrocytoma IDH-mutant, astrocytoma IDH-wildtype, glioblastoma IDH-wildtype, and
        oligodendroglioma IDH-mutant and 1p/19q-codeleted;
  \item LGG (grade 2--3) versus HGG (grade 4);
  \item IDH status, in the late-fusion experiments only.
\end{itemize}

\begin{table}[t]
\centering
\caption{Classification cohort and its fixed split. The segmentation split (345/75/75 of 495 patients) contains it:
every classification test patient is a segmentation test patient, and the same holds for validation and training.}
\label{tab:cohort}
\small
\begin{tabular}{lrrrr}
\toprule
 & Train & Validation & Test & Total \\
\midrule
Grade 2 & 38 & 9 & 9 & 56 \\
Grade 3 & 31 & 6 & 6 & 43 \\
Grade 4 & 43 & 10 & 10 & 63 \\
\midrule
Astrocytoma, IDH-mutant & 47 & 11 & 10 & 68 \\
Astrocytoma, IDH-wildtype & 17 & 3 & 4 & 24 \\
Glioblastoma, IDH-wildtype & 39 & 9 & 9 & 57 \\
Oligodendroglioma, IDH-mutant, 1p/19q-codeleted & 9 & 2 & 2 & 13 \\
\midrule
Patients & 112 & 25 & 25 & 162 \\
\bottomrule
\end{tabular}
\end{table}

\subsection{Splits and leakage controls}
One patient-level split was created before any training and used by every stage. The 162-patient cohort was split
70/15/15 with stratification on grade and diagnosis (112/25/25). The 495 segmentation patients were split 345/75/75
so that classification test patients are segmentation test patients, and likewise for validation and training.
Self-supervised pretraining used only the segmentation training and validation patients. No model, selector or
scaler saw a classification test patient before the final evaluation.

``Clinical'' input means variables known before surgery: age, sex, and whether a biopsy preceded the MRI.
IDH, 1p/19q, MGMT and extent of resection come from tissue or surgery, and IDH and 1p/19q define the WHO~2021
diagnosis. They were excluded from all main models. They appear only in rows labelled \emph{leaky reference}, which
add them to the inputs (the \emph{full clinical vector}) to quantify how much such variables inflate performance.

\subsection{Preprocessing}
Volumes were reoriented to RAS, resampled to 1~mm isotropic spacing and z-score normalised per volume over non-zero
voxels. Networks worked on $96^3$ patches. Metrics were computed on the 1~mm grid.

\subsection{Segmentation}
\label{sec:seg}
We compared:
\begin{itemize}
  \item SwinUNETR~\cite{hatamizadeh2022swinunetr}, trained from scratch and initialised from self-supervised
        pretraining;
  \item SegResNet~\cite{myronenko2019segresnet};
  \item a 3D residual U-Net~\cite{cicek2016unet3d} with channels 32--320, and, as a configuration sensitivity
        analysis, the same U-Net in a narrower configuration (channels 16--256, dropout 0.1);
  \item SegMamba~\cite{xing2024segmamba,gu2023mamba}.
\end{itemize}
All were implemented in MONAI~\cite{cardoso2022monai} and trained with one budget:
\begin{itemize}
  \item Dice plus cross-entropy loss;
  \item AdamW~\cite{loshchilov2019adamw} with learning rate $10^{-4}$ and weight decay $10^{-5}$, 3 warm-up epochs then
        cosine decay, for 100 epochs;
  \item batches of 2 volumes with 2 crops each, sampled 1:1 around tumour and background voxels;
  \item mixed precision and gradient clipping at norm 1;
  \item random flips, small affine transforms, Gaussian noise, intensity shift and scaling, and bias fields as
        augmentation.
\end{itemize}
The checkpoint with the best validation Dice, evaluated every 5 epochs, was kept.

Under the shared learning rate, SegMamba's validation Dice peaked at epoch 35 and then declined while its training
loss kept falling. We therefore added one sensitivity run at half the learning rate ($5 \times 10^{-5}$). It was
added after the main segmentation results were known, is reported as its own row, and is never the reference model
in the statistical comparison.

Inference used sliding windows ($96^3$, 50\% overlap, Gaussian weighting) and a 0.5 threshold. We report Dice, IoU,
the 95th-percentile Hausdorff distance (HD95), the normalised surface distance at 1~mm
(NSD)~\cite{nikolov2021nsd}, sensitivity and precision, on the raw predictions of every
model~\cite{maierhein2024metrics}. Confidence intervals are 95\% percentile bootstrap intervals over patients
(2,000 resamples)~\cite{efron1993bootstrap}. Every model was compared with the best model by Dice using two-sided
paired Wilcoxon signed-rank tests~\cite{wilcoxon1945} on Dice and HD95, with Holm correction~\cite{holm1979}.

\paragraph{Self-supervised pretraining.} The SwinUNETR encoder was pretrained for 50 epochs with four
objectives~\cite{tang2022ssl}:
\begin{itemize}
  \item reconstruction of masked $16^3$ cubes (60\% masked, weight 1.0)~\cite{he2022mae};
  \item prediction of a $k \times 90^\circ$ rotation (weight 0.5);
  \item an InfoNCE contrastive loss between two augmented views (temperature 0.5, weight
        1.0)~\cite{oord2018cpc,chen2020simclr};
  \item an unmasked reconstruction (weight 0.1).
\end{itemize}
We used AdamW (learning rate $10^{-4}$) with gradient accumulation to an effective batch of 8.

\subsection{Shared backbone, predicted masks and features}
All classification stages share one segmentation backbone. We used SegMamba to test whether features from an
encoder with long-range, state-space context add to radiomics. The backbone was fixed after the segmentation
benchmark but before any classifier was trained, and it is not the best segmenter in Table~\ref{tab:seg}. Of its
two runs, the one with the better validation Dice (learning rate $5 \times 10^{-5}$; 0.893 against 0.883) was
chosen before its test result was known.

Its predicted masks, reduced to their largest connected component, centre the evaluation crops and define the
radiomics region for every classification patient. This cleanup step was chosen from validation HD95 before any
classifier was trained.

\paragraph{Radiomics.} 102 original-image features (first order, shape, GLCM, GLRLM, GLSZM, GLDM) were extracted with
PyRadiomics~\cite{vangriethuysen2017pyradiomics} from the predicted masks (bin width 5, intensity normalisation,
1~mm resampling). Thirty were kept by recursive feature elimination with a 50-tree random forest, fit on training data
only.

\paragraph{Deep features.} For late fusion we pooled SegMamba's encoder stages 2, 3 and 4 and its decoder features
(192 + 384 + 768 + 48 = 1,392 values) with a masked mean inside the predicted tumour probability. The backbone was
trained with tumour masks only, never with class labels, so these features can be cross-validated over the
classification cohort without label leakage. It has, however, seen the images and masks of the classification
training and validation patients during segmentation training.

\subsection{End-to-end classifiers}
\label{sec:e2e}
The end-to-end model is a multi-task network on the SegMamba backbone, initialised from the segmentation model.
\begin{itemize}
  \item \emph{Segmentation head.} It stays trained (Dice plus cross-entropy, weight 1).
  \item \emph{Tokens.} The stage-4 bottleneck and the decoder features are pooled with a masked mean weighted by the
        head's own tumour probability. They are projected to 128-dimensional tokens, together with one token per
        tabular input (clinical, radiomics).
  \item \emph{Fusion.} The tokens are fused with learnable task tokens by a two-layer, four-head
        transformer~\cite{vaswani2017attention}.
  \item \emph{Training.} Random $96^3$ crops containing tumour, modality dropout (0.3), auxiliary per-token heads
        (weight 0.3), and a focal loss ($\gamma = 2$)~\cite{lin2017focal} with balanced class weights (task weights:
        grade 1.0, diagnosis 0.5).
  \item \emph{Optimisation.} AdamW with learning rates $2 \times 10^{-4}$ for the heads and $10^{-5}$ for the
        backbone, weight decay 0.05, 5 warm-up epochs, cosine decay, up to 100 epochs. The backbone was frozen for
        the first 10 epochs.
  \item \emph{Checkpoint.} Chosen on the exponentially smoothed (factor 0.6) validation negative log-likelihood of
        the classification heads, with patience 25.
  \item \emph{Inference.} Validation and test predictions average the probabilities over the 8 flips of the crop.
        At test time the crop is centred on the predicted tumour.
\end{itemize}

The main comparison is a ladder of the same model in which only the inputs change: image only, image + clinical,
and image + clinical + radiomics. The leaky reference adds the full clinical vector to the last two rungs. Four
alternative designs complete the set:
\begin{itemize}
  \item \emph{cascade}: frozen backbone, whole-crop pooling, concatenation fusion;
  \item \emph{multi-scale}: encoder stages 2--4 plus decoder, hard tumour mask;
  \item \emph{Trinity}: cross-attention fusion over stages 1--3 with a supervised contrastive term~\cite{khosla2020supcon};
  \item \emph{SegMamba-v3}: bottleneck pooled in the tumour core and in a peritumoral ring (5-voxel dilation), with
        a dedicated LGG/HGG head.
\end{itemize}
The checkpoint selection, the test-time flips and the initial backbone freeze were fixed before any classifier was
trained.

For models without an LGG/HGG head, LGG versus HGG is derived from the grade head. $P(\text{grade 4})$ is the score
for the AUC, and the call is HGG exactly when the grade head predicts grade 4. \emph{This rule was changed after the
first test evaluation.} Originally the call was HGG only when $P(\text{grade 4}) > 0.5$. Under that rule three models
called every test patient LGG despite LGG/HGG AUCs of 0.76--0.89, and a model could call a patient grade 4 and LGG at
the same time. The change affects only LGG/HGG balanced accuracy; no AUC changed and nothing was retrained. We
therefore treat the LGG/HGG AUC as the primary LGG/HGG metric, and the first evaluation's balanced accuracies are in
the supplement.

\subsection{Late fusion}
Late fusion used repeated stratified 5-fold CV (10 repeats, 50 folds) over all 162 patients. Every imputer, scaler,
selector and classifier was refit inside each training fold: RFE to 30 radiomics features, and an ANOVA top-32 of the
deep features. The classifiers were:
\begin{itemize}
  \item logistic regression ($C = 0.5$);
  \item a random forest;
  \item a one-hidden-layer MLP (64 units);
  \item XGBoost~\cite{chen2016xgboost} (200 trees, depth 4, learning rate 0.05, balanced sample weights).
\end{itemize}
All were implemented with scikit-learn~\cite{pedregosa2011sklearn}, with hyperparameters fixed before any result. A
majority-class floor was always included.

The feature sets mirror the ladder: age + sex, clinical, radiomics, clinical + radiomics, image (deep features),
image + clinical, image + clinical + radiomics, and four leaky references. XGBoost is the pre-specified primary model,
and logistic regression is reported alongside. As a fully clean confirmation, every model was also fit once on
training + validation and scored on the 25 test patients. XGBoost feature attributions were computed with
TreeSHAP~\cite{lundberg2017shap,lundberg2020tree}.

Late fusion ran on CPU (scikit-learn 1.2.1, XGBoost 2.1.3). An earlier run of this pipeline on the same folds under
other library versions differed in CV macro AUC by at most 0.009 for XGBoost, logistic regression and the random
forest, and by at most 0.028 for the MLP.

\subsection{Explanations and classification metrics}
For the first seed of each end-to-end classifier we recorded:
\begin{itemize}
  \item the attention of the task tokens over the input tokens;
  \item the drop in true-class probability when each token is zeroed (token occlusion);
  \item Grad-CAM~\cite{selvaraju2017gradcam} maps of the pooled feature maps on the unflipped crop.
\end{itemize}
Classification is reported as balanced accuracy, macro F1 and macro one-versus-rest AUC, with 95\% patient-level
bootstrap CIs (2,000 resamples; the hard calls are resampled with the probabilities).

\subsection{Statistical comparisons}
\label{sec:stats}
The comparisons below were fixed before any of them was tested, and all are reported.
\begin{itemize}
  \item \emph{End-to-end models, on the test patients.} Each rung of the ladder was compared with the rung below
        (clinical token, radiomics token), and each alternative design with the ladder model that has the same
        inputs. Differences in AUC were tested with DeLong's test~\cite{delong1988} for LGG versus HGG and with a
        paired bootstrap of the macro AUC (2,000 patient resamples) for grade and diagnosis, with Holm correction
        within each task.
  \item \emph{Late fusion, over the CV folds.} For XGBoost and logistic regression we compared clinical + radiomics
        with clinical, image + clinical with clinical, image + clinical + radiomics with clinical + radiomics, and
        radiomics with image. The 50 folds are identical for every feature set and model of a task, so the fold
        scores pair up. Because folds share most of their training data, we used the corrected resampled
        $t$-test~\cite{nadeau2003inference,bouckaert2004replicability}, which inflates the variance of the mean
        fold difference by the test/train size ratio, with Holm correction within each task and model.
\end{itemize}

\paragraph{Ground-truth-mask analysis.} To measure what segmentation errors cost, the same trained models were
scored with the expert masks in place of the predicted ones. For the end-to-end classifiers the masks set the crop
centre and the radiomics; inside the network the tumour weighting remains the model's own segmentation head. For
late fusion, the radiomics and deep features were recomputed from the expert masks, and the fold pipelines were
refit on the same folds. Nothing was retrained or reselected, and the same tests compare the two mask sources.

\subsection{Further analyses}
\begin{itemize}
  \item \emph{Segmentation against tumour size.} Spearman correlation between per-patient Dice and ground-truth
        volume, and Dice per volume tertile of the 75 test patients.
  \item \emph{Calibration and clinical usefulness.} Multi-class Brier score~\cite{brier1950} and expected
        calibration error (ECE, 10 bins)~\cite{guo2017calibration}, for the end-to-end models on the test patients
        and for late fusion on out-of-fold predictions of all 162 patients (one 5-fold pass). For LGG versus HGG and
        IDH we also drew reliability diagrams and decision curves~\cite{vickers2006decision}. All models were trained
        with balanced class weights, which shifts probabilities towards the rarer classes by design.
  \item \emph{Learning curves.} Late fusion was fit on stratified 25, 50, 75 and 100\% subsamples of each training
        fold (5-fold CV, 4 repeats) and scored on the full test folds. Every feature set saw the same subsamples.
  \item \emph{Radiomics stability.} Agreement of each radiomics feature between predicted and expert masks, as
        ICC(2,1)~\cite{shrout1979icc,koo2016icc} and Lin's concordance correlation~\cite{lin1989ccc}. Only the 50
        classification validation and test patients count, because the backbone's predicted masks are nearly
        perfect on its own training patients.
  \item \emph{Subgroups.} Late-fusion macro AUC (XGBoost and logistic regression) by sex and by age band (under 40,
        40--59, 60 and over) on the out-of-fold predictions, with patient-level bootstrap CIs. No AUC is given where a
        class has fewer than five patients in the subgroup. Age varies less within a band, so these also show what
        imaging contributes when age has less to offer. For the binary tasks we also scored patients by age directly,
        with no model, as a reference.
\end{itemize}

\subsection{Post hoc checks}
\label{sec:posthoc}
The two checks below were added after the results above were known, and are exploratory.
\begin{itemize}
  \item \emph{Validation--test gap.} Is a model's gap between validation and test AUC larger than chance assignment of
        patients to the two splits would give? Each model's per-patient probabilities were held fixed and the 50
        patients were re-assigned at random to two groups, within each class, so that the group sizes and class mix
        stayed the same (10,000 re-assignments, two-sided permutation $p$~\cite{phipson2010permutation}, Holm
        correction within each task and family of models). The checkpoints were selected on the validation patients,
        which if anything favours validation. For a like-for-like reference, late fusion was fit on the same 112
        training patients as the end-to-end models. The same re-assignments served every model of a task, so each
        end-to-end model of the main ladder could also be compared with late fusion on the same inputs (Holm within
        each task).
  \item \emph{Shuffled-label control.} On shuffled labels no real association exists, so the AUC must be 0.5 unless
        information leaks~\cite{ojala2010permutation}. We ran 100 permutations for diagnosis, LGG versus HGG and IDH
        status, with XGBoost and logistic regression on radiomics (102 features, recursive elimination to 30) and on
        deep features (1,392 features, ANOVA to 32), in one 5-fold pass per permutation. Each was fit two ways: with
        the imputer, scaler and selector fit inside each training fold, as everywhere in this study, and with them fit
        once on all patients before the folds, the common mistake. The real labels went through the same two
        procedures on the 50 folds of the main CV.
\end{itemize}

% ---------------------------------------------------------------------------------------------------------
\section{Results}
All end-to-end and segmentation results are from one seed. Late-fusion results are means over the 50 CV folds.

\subsection{Segmentation}
\begin{table}[t]
\centering
\caption{Whole-tumour segmentation on the 75 test patients (one seed). Dice and HD95 with 95\% bootstrap CIs.
$p$: paired Wilcoxon against SegResNet, Holm-corrected. Train: hours on one A100. Inference: seconds per volume.
$^\dagger$Sensitivity run added after the main results (Section~\ref{sec:seg}). $^\ast$Training time partly
estimated from log timestamps.}
\label{tab:seg}
\small
\resizebox{\textwidth}{!}{%
\begin{tabular}{lcccccccc}
\toprule
Model & Dice & HD95 (mm) & NSD & $p_\text{Dice}$ & $p_\text{HD95}$ & Params (M) & Train (h) & Infer. (s) \\
\midrule
SegResNet & \textbf{0.865}\ci{0.828}{0.899} & \textbf{11.7}\ci{7.5}{16.4} & \textbf{0.715} & -- & -- & 4.7 & 0.9 & 0.57 \\
SegMamba, lr $5\times10^{-5}$ $^\dagger$ & 0.847\ci{0.797}{0.890} & 22.3\ci{12.8}{32.8} & 0.704 & 0.84 & 0.007 & 67.4 & 3.2 & 2.09 \\
SwinUNETR & 0.843\ci{0.791}{0.886} & 23.3\ci{15.5}{31.3} & 0.686 & 0.22 & 0.002 & 62.2 & 2.7$^\ast$ & 1.55 \\
SwinUNETR + SSL & 0.842\ci{0.797}{0.883} & 20.1\ci{12.3}{28.2} & 0.678 & 0.004 & 0.001 & 62.2 & 2.8$^\ast$ & 1.65 \\
3D U-Net & 0.838\ci{0.793}{0.879} & 26.3\ci{18.2}{35.4} & 0.654 & $<10^{-6}$ & $<10^{-6}$ & 12.9 & 1.2 & 0.15 \\
SegMamba & 0.832\ci{0.781}{0.877} & 14.2\ci{9.7}{19.2} & 0.670 & $<10^{-3}$ & 0.005 & 67.4 & 3.2 & 2.09 \\
3D U-Net, narrow config. & 0.728\ci{0.675}{0.781} & 131.2\ci{122.2}{139.3} & 0.406 & $<10^{-12}$ & $<10^{-12}$ & 4.8 & 1.2 & 0.14 \\
\bottomrule
\end{tabular}}
\end{table}

SegResNet had the highest mean Dice and the lowest HD95 (Table~\ref{tab:seg}, Fig.~\ref{fig:seg}). Compared with it:
\begin{itemize}
  \item SegMamba at the lower learning rate and SwinUNETR from scratch were not significantly different in Dice
        ($p = 0.84$ and $0.22$) but had significantly larger HD95;
  \item SegMamba at the shared learning rate was worse on both.
\end{itemize}
Median Dice was 0.91--0.93 for every model except the narrow U-Net configuration (0.815). The means are pulled down
by three to six patients per model with Dice below 0.5. Two patients fail for every model: a 1.6~mL tumour that no
model detects, and a glioblastoma on which every model labels mostly the wrong tissue (precision 0.04--0.18).

Self-supervised pretraining did not help. SwinUNETR + SSL against SwinUNETR from scratch differed by $-0.001$ Dice
($p = 0.08$, unadjusted) and $-3.2$~mm HD95 ($p = 0.27$). U-Net results depended strongly on configuration: under the
same budget the wider U-Net gained 0.110 Dice and lost 105~mm HD95 against the narrow one ($p < 10^{-12}$ for both).
The narrow configuration's HD95 of 131~mm comes from scattered false positives across the brain (Fig.~\ref{fig:seg}b).
SegResNet was also the smallest and among the fastest models (4.7~M parameters and 0.57~s per volume, against
67.4~M and 2.09~s for SegMamba).

Failures concentrated in small tumours (Fig.~\ref{fig:volume}). Dice rose with tumour volume for every model
(Spearman $\rho$ 0.45--0.56, and 0.74 for the narrow U-Net; all $p < 10^{-4}$). In the smallest third of the test
tumours (below 50~mL), median Dice was 0.81--0.84 (narrow U-Net 0.62), against 0.93--0.94 in the largest third
(above 103~mL). Of the 40 failures with Dice below 0.5 across the seven models, 38 were in the smallest third.

\begin{figure}[t]
\centering
\includegraphics[width=\textwidth]{segmentation_dice_vs_volume.png}
\caption{Per-patient test Dice against ground-truth tumour volume (log scale), one panel per model. Dashed line:
Dice 0.5.}
\label{fig:volume}
\end{figure}

\begin{figure}[t]
\centering
\begin{subfigure}{\textwidth}\centering
  \includegraphics[width=\textwidth]{segmentation_dice.png}
  \caption{Per-patient test Dice (green triangle: mean).}
\end{subfigure}
\begin{subfigure}{\textwidth}\centering
  \includegraphics[width=\textwidth]{segmentation_cases.png}
  \caption{Best, median and worst test case of SegResNet, for every model (green: ground truth, red: prediction).}
\end{subfigure}
\caption{Whole-tumour segmentation on the 75 test patients.}
\label{fig:seg}
\end{figure}

\subsection{What leaky inputs would explain}
With the full clinical vector and no imaging at all, XGBoost reached a CV diagnosis AUC of $0.930 \pm 0.022$
(logistic regression $0.949 \pm 0.023$) and IDH AUC 1.000 (Table~\ref{tab:kfold}). End to end, the leaky reference
with radiomics reached a test diagnosis balanced accuracy of 0.944 (0.857--1.000) and AUC 0.975
(Table~\ref{tab:cls}). Such numbers therefore need no imaging, and every comparison below excludes these variables.

\subsection{End-to-end classification}
\begin{table}[t]
\centering
\caption{End-to-end classifiers on the 25 test patients (one seed). Balanced accuracy (BA) and macro AUC with 95\%
bootstrap CIs. Chance: BA 0.333 (grade) and 0.250 (diagnosis); AUC 0.5. LGG/HGG for all models except SegMamba-v3
is derived from the grade head (Section~\ref{sec:e2e}).}
\label{tab:cls}
\small
\resizebox{\textwidth}{!}{%
\begin{tabular}{lccccc}
\toprule
 & \multicolumn{2}{c}{Grade (3-way)} & \multicolumn{2}{c}{Diagnosis (4-way)} & LGG vs HGG \\
\cmidrule(lr){2-3}\cmidrule(lr){4-5}\cmidrule(lr){6-6}
Model & BA & AUC & BA & AUC & AUC \\
\midrule
\multicolumn{6}{l}{\emph{Main ladder (same model, inputs change)}} \\
Image only & 0.530\ci{0.40}{0.64} & 0.761\ci{0.63}{0.87} & 0.364\ci{0.27}{0.54} & 0.799\ci{0.66}{0.93} & 0.887\ci{0.72}{1.00} \\
Image + clinical & 0.456\ci{0.31}{0.60} & 0.751\ci{0.58}{0.90} & 0.457\ci{0.32}{0.70} & 0.869\ci{0.74}{0.98} & 0.853\ci{0.66}{1.00} \\
Image + clinical + radiomics & 0.548\ci{0.34}{0.75} & 0.808\ci{0.65}{0.93} & 0.744\ci{0.55}{0.94} & 0.916\ci{0.84}{0.99} & 0.867\ci{0.68}{1.00} \\
\midrule
\multicolumn{6}{l}{\emph{Alternative designs (image + clinical, + radiomics where marked)}} \\
Cascade & 0.493\ci{0.35}{0.63} & 0.674\ci{0.53}{0.81} & 0.419\ci{0.33}{0.62} & 0.857\ci{0.77}{0.95} & 0.793\ci{0.59}{0.96} \\
Multi-scale & 0.563\ci{0.44}{0.67} & 0.708\ci{0.54}{0.86} & 0.394\ci{0.30}{0.59} & 0.804\ci{0.65}{0.95} & 0.873\ci{0.69}{1.00} \\
Trinity (+ radiomics) & 0.526\ci{0.40}{0.64} & 0.739\ci{0.56}{0.89} & 0.369\ci{0.26}{0.56} & 0.788\ci{0.62}{0.94} & 0.760\ci{0.50}{0.97} \\
SegMamba-v3 (+ radiomics) & 0.567\ci{0.38}{0.77} & 0.766\ci{0.59}{0.91} & 0.757\ci{0.55}{0.90} & 0.892\ci{0.80}{0.98} & 0.840\ci{0.63}{1.00} \\
\midrule
\multicolumn{6}{l}{\emph{Leaky reference (uses post-operative molecular and surgical data)}} \\
Image + full clinical & 0.563\ci{0.37}{0.76} & 0.841\ci{0.72}{0.95} & 0.882\ci{0.71}{1.00} & 0.969\ci{0.91}{1.00} & 0.947\ci{0.84}{1.00} \\
Image + full clinical + radiomics & 0.622\ci{0.42}{0.82} & 0.870\ci{0.74}{0.97} & 0.944\ci{0.86}{1.00} & 0.975\ci{0.93}{1.00} & 0.927\ci{0.78}{1.00} \\
\bottomrule
\end{tabular}}
\end{table}

On the test split, adding radiomics to image + clinical raised diagnosis balanced accuracy from 0.457 to 0.744 and AUC
from 0.869 to 0.916 (Table~\ref{tab:cls}). The AUC difference was not significant: $+0.047$ (95\% CI $-0.012$ to
$0.113$; $p = 0.14$, Holm-adjusted 0.57). Adding the clinical token gave $+0.070$ for diagnosis ($-0.010$ to
$0.152$; $p = 0.10$). Among the pre-specified comparisons (Section~\ref{sec:stats}), none was significant after Holm
correction for any task. The largest difference was the Trinity design's deficit against the ladder model with the
same inputs (diagnosis $-0.127$, $-0.238$ to $-0.017$; $p = 0.019$, adjusted 0.11).
\todo{supplement table: \texttt{paper\_tables/classification\_tests}.}

The weaker models' balanced accuracies are low despite reasonable AUCs because their calls collapse onto the two
large classes (Fig.~\ref{fig:cm}):
\begin{itemize}
  \item Image only, image + clinical, cascade, multi-scale and Trinity predicted grade 3 at most twice, astrocytoma
        IDH-wildtype at most once, and oligodendroglioma never.
  \item Image only had a mean top-class probability of 0.36 for diagnosis (uniform: 0.25), the same for correct and
        incorrect calls.
  \item The image + clinical + radiomics model predicted all four classes: astrocytoma IDH-wildtype 4/4,
        glioblastoma 7/9, astrocytoma IDH-mutant 7/10, oligodendroglioma 1/2.
  \item SegMamba-v3's balanced accuracy is inflated by over-predicting the rare oligodendroglioma class (9 calls
        for 2 patients; accuracy 0.68).
\end{itemize}
None of the alternative designs clearly improved on the plain ladder; all differences lie well inside the confidence
intervals.

The probabilities were poorly calibrated, as expected with balanced class weights. Diagnosis Brier scores were 0.65
for image only (uniform probabilities score 0.75), 0.46 for image + clinical + radiomics and 0.24 for the leaky
reference with radiomics, with ECEs of 0.20--0.31. \todo{supplement table: \texttt{paper\_tables/calibration\_end\_to\_end}.}

\begin{figure}[t]
\centering
\begin{subfigure}{0.32\textwidth}\centering
  \includegraphics[width=\textwidth]{cls_image_only_diagnosis_confusion.png}
  \caption{Image only.}
\end{subfigure}
\begin{subfigure}{0.32\textwidth}\centering
  \includegraphics[width=\textwidth]{cls_image_clinical_radiomics_diagnosis_confusion.png}
  \caption{Image + clinical + radiomics.}
\end{subfigure}
\begin{subfigure}{0.32\textwidth}\centering
  \includegraphics[width=\textwidth]{cls_leaky_reference_radiomics_diagnosis_confusion.png}
  \caption{Leaky reference + radiomics.}
\end{subfigure}
\caption{Test confusion matrices for four-way diagnosis (rows: truth). The image-only model collapses onto
astrocytoma IDH-mutant and glioblastoma; the leaky reference is nearly perfect because its inputs encode the label.}
\label{fig:cm}
\end{figure}

\subsection{Late fusion and what imaging adds}
\begin{table}[t]
\centering
\caption{Late fusion, repeated stratified 5-fold CV over the 162 patients (mean $\pm$ SD of macro AUC over 50 folds),
XGBoost. Majority class: 0.500 for every task. ``Image'' = label-free deep features of the segmentation backbone.}
\label{tab:kfold}
\small
\begin{tabular}{lcccc}
\toprule
Features & Grade & Diagnosis & LGG vs HGG & IDH \\
\midrule
Age + sex & $0.569 \pm 0.047$ & $0.649 \pm 0.044$ & $0.656 \pm 0.082$ & $0.846 \pm 0.058$ \\
Clinical & $0.564 \pm 0.049$ & $0.652 \pm 0.045$ & $0.653 \pm 0.080$ & $0.839 \pm 0.057$ \\
Radiomics & $0.736 \pm 0.064$ & $0.790 \pm 0.062$ & $0.795 \pm 0.079$ & $0.823 \pm 0.063$ \\
Clinical + radiomics & $\mathbf{0.762 \pm 0.059}$ & $0.830 \pm 0.051$ & $\mathbf{0.829 \pm 0.063}$ & $0.899 \pm 0.045$ \\
Image & $0.692 \pm 0.058$ & $0.751 \pm 0.048$ & $0.786 \pm 0.085$ & $0.822 \pm 0.069$ \\
Image + clinical & $0.704 \pm 0.059$ & $0.803 \pm 0.048$ & $0.798 \pm 0.073$ & $\mathbf{0.907 \pm 0.048}$ \\
Image + clinical + radiomics & $0.748 \pm 0.051$ & $\mathbf{0.831 \pm 0.046}$ & $0.819 \pm 0.076$ & $0.906 \pm 0.048$ \\
\midrule
\multicolumn{5}{l}{\emph{Leaky reference (post-operative molecular and surgical data)}} \\
Full clinical vector & $0.737 \pm 0.056$ & $0.930 \pm 0.022$ & $0.862 \pm 0.035$ & $1.000 \pm 0.000$ \\
Molecular + surgical & $0.740 \pm 0.057$ & $0.944 \pm 0.025$ & $0.905 \pm 0.037$ & $1.000 \pm 0.000$ \\
Image + full clinical & $0.796 \pm 0.051$ & $0.944 \pm 0.023$ & $0.905 \pm 0.044$ & $1.000 \pm 0.000$ \\
Image + full clinical + radiomics & $0.811 \pm 0.043$ & $0.956 \pm 0.019$ & $0.910 \pm 0.046$ & $1.000 \pm 0.000$ \\
\bottomrule
\end{tabular}
\end{table}

\begin{table}[t]
\centering
\caption{Pre-specified late-fusion comparisons over the 50 CV folds: difference in macro AUC, with the corrected
resampled $t$-test $p$ in brackets, Holm-adjusted within each task and model. ``Deep'' = deep features.}
\label{tab:kfoldtests}
\small
\begin{tabular}{lcccc}
\toprule
Comparison & Grade & Diagnosis & LGG vs HGG & IDH \\
\midrule
\multicolumn{5}{l}{\emph{XGBoost}} \\
Radiomics added to clinical & $+0.198$ ($<0.001$) & $+0.178$ ($<0.001$) & $+0.176$ (0.004) & $+0.060$ (0.0502) \\
Deep added to clinical & $+0.140$ ($<0.001$) & $+0.150$ ($<0.001$) & $+0.145$ (0.020) & $+0.068$ (0.058) \\
Deep added to clinical + radiomics & $-0.014$ (0.54) & $+0.000$ (0.99) & $-0.010$ (1) & $+0.007$ (1) \\
Radiomics against deep & $+0.044$ (0.54) & $+0.039$ (0.58) & $+0.008$ (1) & $+0.001$ (1) \\
\midrule
\multicolumn{5}{l}{\emph{Logistic regression}} \\
Radiomics added to clinical & $+0.074$ (0.21) & $+0.103$ (0.004) & $+0.059$ (0.79) & $-0.005$ (1) \\
Deep added to clinical & $+0.043$ (0.69) & $+0.076$ (0.046) & $+0.048$ (0.79) & $+0.015$ (1) \\
Deep added to clinical + radiomics & $-0.017$ (0.69) & $-0.023$ (0.42) & $-0.015$ (1) & $+0.008$ (1) \\
Radiomics against deep & $+0.040$ (0.69) & $+0.043$ (0.42) & $+0.015$ (1) & $-0.001$ (1) \\
\bottomrule
\end{tabular}
\end{table}

Repeated CV gives the more reliable comparison (Tables~\ref{tab:kfold} and~\ref{tab:kfoldtests}). For diagnosis,
imaging added to clinical variables with both classifiers and both encodings: radiomics by 0.178 (XGBoost) and 0.103
(logistic regression), deep features by 0.150 and 0.076, all significant after correction. For grade and LGG/HGG the
gain depended on the classifier. It was significant for XGBoost and not for logistic regression, largely because
XGBoost handled the clinical variables alone much worse (grade 0.564 against 0.648, LGG/HGG 0.653 against 0.746). The
highest AUCs for these two tasks still came from XGBoost (0.762 and 0.829, against 0.722 and 0.804).

The two encodings of the images overlapped. Adding deep features to clinical + radiomics gained nothing (XGBoost
diagnosis $0.831$ against $0.830$; logistic regression $0.807$ against $0.830$). Radiomics alone scored higher than deep
features alone for grade, diagnosis and LGG/HGG, but none of these differences was significant.

For IDH status, age and sex alone reached $0.846$ (XGBoost) and $0.895$ (logistic regression). The best imaging models,
image + clinical with either classifier, reached 0.907 and 0.906, which is 0.06 and 0.01 above age and sex. No gain
over clinical variables survived correction (Table~\ref{tab:kfoldtests}).

\paragraph{Learning curves.} More training patients did not help age + sex. From about 32 to 130 training patients, its
diagnosis AUC stayed at 0.725--0.713 with logistic regression and 0.664--0.645 with XGBoost (Fig.~\ref{fig:learning}).
Imaging did gain: radiomics rose from 0.702 to 0.783 (XGBoost) and from 0.722 to 0.798 (logistic regression), and deep
features from 0.651 and 0.679 to 0.749 with both. At about 32 patients, age + sex by logistic regression still matched
radiomics for diagnosis (0.725 against 0.722), and radiomics passed it only from about 65. For IDH, the same age + sex
model (0.895--0.897 at every size) stayed above radiomics or deep features alone at every size (at most 0.825).

Deep features did not clearly catch up with radiomics as the training set grew. With XGBoost, radiomics led by
0.051, 0.039, 0.046 and 0.034 for diagnosis and by 0.023, 0.053, 0.070 and 0.040 for grade across the four sizes. For
LGG/HGG and IDH the lead was gone at the full size ($-0.001$ and $-0.003$). The SDs over the test folds, 0.05--0.10 at
the full size, are larger than any of these differences. The XGBoost diagnosis curves flatten between about 97 and 130
patients, whereas logistic regression was still rising, so we cannot say that 130 patients exhaust what T2 offers.

\begin{figure}[t]
\centering
\begin{subfigure}{0.49\textwidth}\centering
  \includegraphics[width=\textwidth]{learning_curve_diagnosis.png}
  \caption{WHO~2021 diagnosis.}
\end{subfigure}
\begin{subfigure}{0.49\textwidth}\centering
  \includegraphics[width=\textwidth]{learning_curve_idh.png}
  \caption{IDH status.}
\end{subfigure}
\caption{Late-fusion macro AUC against training-set size (XGBoost; mean $\pm$ SD over 20 test folds). Every feature set
saw the same subsamples, and the test folds were not subsampled.}
\label{fig:learning}
\end{figure}

\paragraph{Calibration and decision curves.} On the out-of-fold predictions the Brier score of late fusion
(multi-class form, so a constant at the cohort's prevalence scores 0.48 for LGG/HGG and 0.50 for IDH) improved with
imaging for LGG/HGG: age + sex 0.526 (XGBoost) and 0.407 (logistic regression), clinical + radiomics 0.338 and 0.341.
For IDH it did not move much: age + sex 0.291 and 0.265, image + clinical + radiomics 0.262 and 0.266. ECE was 0.07--0.17
across the five feature sets drawn in Fig.~\ref{fig:decision}. The decision curves show the same picture
(Fig.~\ref{fig:decision}). For LGG/HGG, the models with clinical + radiomics gave the highest net benefit above a
threshold of about 0.25, ahead of age + sex with either classifier. For IDH, every XGBoost model beat treating everyone
above a threshold of about 0.2, but imaging gave little: image + clinical + radiomics led age + sex (logistic regression)
by at most 0.02 below a threshold of 0.45 and fell behind it above 0.5, and radiomics or deep features alone stayed below
it from 0.2. These curves inherit the cohort's designed prevalence (39\% grade 4, 50\% IDH-wildtype), so they compare the
models with each other and do not predict clinical benefit.

\begin{figure}[t]
\centering
\begin{subfigure}{0.49\textwidth}\centering
  \includegraphics[width=\textwidth]{decision_curve_lgg_hgg.png}
  \caption{LGG versus HGG.}
\end{subfigure}
\begin{subfigure}{0.49\textwidth}\centering
  \includegraphics[width=\textwidth]{decision_curve_idh.png}
  \caption{IDH status.}
\end{subfigure}
\caption{Decision curves of late fusion (XGBoost, plus age + sex by logistic regression as a no-imaging reference;
out-of-fold predictions of all 162 patients). The prevalence is that of the balanced cohort.}
\label{fig:decision}
\end{figure}

The one-shot hold-out on the 25 test patients agreed within its wide intervals:
\begin{itemize}
  \item XGBoost diagnosis AUC: clinical + radiomics 0.790 [0.63, 0.97], image + clinical + radiomics 0.861 [0.74, 0.96];
  \item IDH AUC: age + sex 0.926 [0.76, 1.00].
\end{itemize}
\todo{full hold-out table in the supplement.}

TreeSHAP on the image + clinical + radiomics XGBoost ranked age first for diagnosis, ahead of any deep or radiomics
feature (Fig.~\ref{fig:shap}). The leading radiomics features were the 10th-percentile intensity and GLCM sum
entropy, a measure of texture heterogeneity. The 32 selected deep features and the 30 selected radiomics features
contributed similarly, in total (1.78 and 1.61) and per feature (0.056 and 0.054), against 0.35 for age alone. XGBoost
thus uses the deep features when they are offered, yet CV shows no gain from adding them, so they appear to be
largely redundant with the radiomics rather than complementary.

\begin{figure}[t]
\centering
\includegraphics[width=0.6\textwidth]{kfold_shap_diagnosis_image_clinical_radiomics_top15.png}
\caption{Top-15 TreeSHAP features for four-way diagnosis (XGBoost, image + clinical + radiomics, fit on training +
validation). \texttt{num\_\_}: clinical, \texttt{deep\_\_}: deep features, \texttt{rad\_\_}: radiomics.}
\label{fig:shap}
\end{figure}

\subsection{Subgroups: does imaging add anything beyond age?}
\begin{table}[t]
\centering
\caption{Late-fusion macro AUC by age band and sex (out-of-fold predictions, 95\% bootstrap CIs). LR: logistic regression.
``Image'': deep features. ``Age itself'': patients scored by age directly, with no model (binary tasks only). A dash marks
a subgroup in which a class has fewer than five patients. Descriptive only: no comparison was tested.}
\label{tab:subgroups}
\small
\resizebox{\textwidth}{!}{%
\begin{tabular}{lccccc}
\toprule
 & Under 40 ($n=53$) & 40--59 ($n=62$) & 60 and over ($n=47$) & Female ($n=63$) & Male ($n=99$) \\
\midrule
\multicolumn{6}{l}{\emph{Grade}} \\
Age + sex, LR & 0.377\ci{0.26}{0.50} & 0.590\ci{0.48}{0.70} & 0.400\ci{0.25}{0.55} & 0.633\ci{0.54}{0.73} & 0.636\ci{0.56}{0.71} \\
Radiomics, XGBoost & 0.630\ci{0.53}{0.74} & 0.727\ci{0.62}{0.83} & 0.638\ci{0.51}{0.76} & 0.742\ci{0.64}{0.83} & 0.671\ci{0.59}{0.74} \\
Radiomics, LR & 0.668\ci{0.54}{0.78} & 0.711\ci{0.60}{0.81} & 0.586\ci{0.45}{0.73} & 0.701\ci{0.59}{0.81} & 0.692\ci{0.62}{0.77} \\
Image, XGBoost & 0.688\ci{0.57}{0.80} & 0.672\ci{0.56}{0.78} & 0.575\ci{0.45}{0.70} & 0.686\ci{0.58}{0.79} & 0.675\ci{0.60}{0.76} \\
Image, LR & 0.601\ci{0.47}{0.72} & 0.640\ci{0.53}{0.74} & 0.595\ci{0.46}{0.74} & 0.651\ci{0.54}{0.76} & 0.642\ci{0.56}{0.73} \\
\midrule
\multicolumn{6}{l}{\emph{Diagnosis}} \\
Age + sex, LR & -- & 0.611\ci{0.52}{0.71} & -- & 0.683\ci{0.59}{0.77} & 0.722\ci{0.66}{0.78} \\
Radiomics, XGBoost & -- & 0.764\ci{0.67}{0.85} & -- & 0.762\ci{0.67}{0.85} & 0.791\ci{0.71}{0.86} \\
Radiomics, LR & -- & 0.774\ci{0.69}{0.86} & -- & 0.793\ci{0.71}{0.87} & 0.795\ci{0.71}{0.87} \\
Image, XGBoost & -- & 0.736\ci{0.64}{0.82} & -- & 0.724\ci{0.63}{0.81} & 0.769\ci{0.69}{0.85} \\
Image, LR & -- & 0.718\ci{0.63}{0.80} & -- & 0.675\ci{0.56}{0.78} & 0.759\ci{0.69}{0.82} \\
\midrule
\multicolumn{6}{l}{\emph{LGG vs HGG}} \\
Age itself & 0.499\ci{0.29}{0.71} & 0.652\ci{0.51}{0.79} & 0.456\ci{0.28}{0.63} & 0.784\ci{0.65}{0.89} & 0.742\ci{0.64}{0.84} \\
Age + sex, LR & 0.453\ci{0.26}{0.65} & 0.618\ci{0.48}{0.76} & 0.481\ci{0.29}{0.67} & 0.768\ci{0.63}{0.88} & 0.721\ci{0.62}{0.82} \\
Radiomics, XGBoost & 0.733\ci{0.56}{0.89} & 0.798\ci{0.67}{0.91} & 0.704\ci{0.54}{0.84} & 0.738\ci{0.59}{0.86} & 0.828\ci{0.74}{0.90} \\
Radiomics, LR & 0.661\ci{0.40}{0.90} & 0.810\ci{0.68}{0.92} & 0.729\ci{0.57}{0.87} & 0.700\ci{0.54}{0.83} & 0.853\ci{0.76}{0.92} \\
Image, XGBoost & 0.786\ci{0.60}{0.95} & 0.713\ci{0.57}{0.84} & 0.729\ci{0.57}{0.86} & 0.690\ci{0.54}{0.82} & 0.810\ci{0.72}{0.89} \\
Image, LR & 0.711\ci{0.50}{0.89} & 0.682\ci{0.53}{0.83} & 0.688\ci{0.53}{0.84} & 0.715\ci{0.56}{0.85} & 0.757\ci{0.66}{0.85} \\
\midrule
\multicolumn{6}{l}{\emph{IDH}} \\
Age itself & -- & 0.762\ci{0.62}{0.88} & 0.569\ci{0.31}{0.81} & 0.868\ci{0.76}{0.95} & 0.912\ci{0.85}{0.96} \\
Age + sex, LR & -- & 0.730\ci{0.58}{0.86} & 0.557\ci{0.29}{0.84} & 0.855\ci{0.74}{0.95} & 0.903\ci{0.84}{0.96} \\
Radiomics, XGBoost & -- & 0.746\ci{0.60}{0.87} & 0.638\ci{0.33}{0.89} & 0.774\ci{0.64}{0.88} & 0.866\ci{0.79}{0.93} \\
Radiomics, LR & -- & 0.715\ci{0.57}{0.86} & 0.524\ci{0.36}{0.69} & 0.733\ci{0.59}{0.85} & 0.827\ci{0.74}{0.90} \\
Image, XGBoost & -- & 0.768\ci{0.65}{0.87} & 0.667\ci{0.40}{0.91} & 0.753\ci{0.63}{0.86} & 0.857\ci{0.78}{0.92} \\
Image, LR & -- & 0.752\ci{0.62}{0.87} & 0.710\ci{0.44}{0.93} & 0.784\ci{0.67}{0.89} & 0.837\ci{0.76}{0.91} \\
\bottomrule
\end{tabular}}
\end{table}

Age dominates many of the results above, so we asked what imaging does when age varies less (Table~\ref{tab:subgroups}).
Age keeps some signal inside a band. Used directly as the score, it reached 0.65 for LGG/HGG and 0.76 for IDH at 40--59,
but only 0.50 and 0.46 for LGG/HGG in the other two bands and 0.57 for IDH at 60 and over. Logistic regression on age +
sex came close to these values. XGBoost on age + sex did not: it fell as low as 0.17 (LGG/HGG at 60 and over), far below
age itself. This is most likely an artefact of a flexible model on one near-discrete feature with little signal on
held-out patients, not a real reverse effect, and it may also contribute to XGBoost's weak clinical baseline in
Table~\ref{tab:kfold}. We therefore use logistic regression as the age + sex reference here.

For grade, diagnosis and LGG/HGG, radiomics and deep features scored above that reference in every band, with both
classifiers. Their confidence intervals excluded 0.5 in 13 of 14 combinations of task, band and encoding for XGBoost and
9 of 14 for logistic regression. For these tasks the T2 signal is therefore not merely a proxy for age. For IDH we could
not show this. At 40--59, radiomics and deep features (XGBoost 0.746 and 0.768, logistic regression 0.715 and 0.752)
were level with age + sex (0.730) and age itself (0.762). At 60 and over there are only five IDH-mutant patients, and
the intervals are too wide to tell.

By sex, the intervals overlapped in every comparison. With imaging alone, point estimates were lower in women for
LGG/HGG and IDH in all eight comparisons (by 0.04--0.15), lower or equal for diagnosis (by 0.00--0.08), and similar or
higher for grade. With 63 women and 99 men, and no test applied, we cannot tell whether these differences are real.

\subsection{Shuffled-label control}
\begin{table}[t]
\centering
\caption{Late fusion on shuffled and real labels. Shuffled: mean $\pm$ SD of macro AUC over 100 permutations, one 5-fold
pass each. Real: the 50 folds of the main CV. ``Inside'': imputer, scaler and selector fit inside each training fold,
as everywhere in this study. ``Before'': fit once on all patients before the folds.}
\label{tab:control}
\small
\begin{tabular}{lllcccc}
\toprule
 & & & \multicolumn{2}{c}{Shuffled labels} & \multicolumn{2}{c}{Real labels} \\
\cmidrule(lr){4-5}\cmidrule(lr){6-7}
Task & Features & Model & Inside & Before & Inside & Before \\
\midrule
Diagnosis & Radiomics & Log. reg. & $0.500 \pm 0.043$ & $0.511 \pm 0.046$ & 0.799 & 0.813 \\
 & & XGBoost & $0.492 \pm 0.047$ & $0.523 \pm 0.051$ & 0.790 & 0.815 \\
 & Deep & Log. reg. & $0.498 \pm 0.048$ & $0.574 \pm 0.043$ & 0.756 & 0.807 \\
 & & XGBoost & $0.497 \pm 0.050$ & $0.570 \pm 0.047$ & 0.751 & 0.794 \\
\midrule
LGG vs HGG & Radiomics & Log. reg. & $0.499 \pm 0.061$ & $0.513 \pm 0.066$ & 0.781 & 0.801 \\
 & & XGBoost & $0.499 \pm 0.069$ & $0.558 \pm 0.071$ & 0.795 & 0.823 \\
 & Deep & Log. reg. & $0.498 \pm 0.055$ & $0.576 \pm 0.053$ & 0.766 & 0.800 \\
 & & XGBoost & $0.506 \pm 0.052$ & $0.575 \pm 0.054$ & 0.786 & 0.820 \\
\midrule
IDH & Radiomics & Log. reg. & $0.504 \pm 0.063$ & $0.517 \pm 0.067$ & 0.816 & 0.832 \\
 & & XGBoost & $0.505 \pm 0.060$ & $0.555 \pm 0.057$ & 0.823 & 0.847 \\
 & Deep & Log. reg. & $0.495 \pm 0.065$ & $0.577 \pm 0.054$ & 0.816 & 0.831 \\
 & & XGBoost & $0.499 \pm 0.061$ & $0.567 \pm 0.055$ & 0.822 & 0.826 \\
\bottomrule
\end{tabular}
\end{table}

On shuffled labels the pipeline as used in this study scored a mean macro AUC of 0.492--0.506 in all 12 configurations
(Table~\ref{tab:control}), so it shows no sign of leakage. Fitting the same steps before the folds gave 0.567--0.577
with the 1,392 deep features and 0.511--0.558 with the 102 radiomics features, on labels with no signal. On real labels
the same mistake inflated macro AUC by 0.004--0.051. That is the size of the differences between feature sets reported
above: radiomics and deep features differ by 0.039 for XGBoost diagnosis. No shuffled run reached the real-label result
in any configuration ($p = 1/101$, the smallest possible with 100 permutations).

\subsection{How reliable is a 25-patient test split?}
\begin{table}[t]
\centering
\caption{Validation against test macro AUC for diagnosis. Both splits have 25 patients and the same class mix. Gap:
test minus validation. Chance: 95\% range of the gap under 10,000 random re-assignments of the 50 patients within
class, with each model's probabilities held fixed. $p$: two-sided permutation test. Adjusted: Holm within each family.
Late fusion was fit on the same 112 training patients as the end-to-end models.}
\label{tab:valtest}
\small
\begin{tabular}{lcccccc}
\toprule
Model & Validation & Test & Gap & Chance & $p$ & Adjusted \\
\midrule
\multicolumn{7}{l}{\emph{End to end}} \\
Image only & 0.381 & 0.799 & $+0.418$ & $-0.252$ to $+0.250$ & 0.0006 & 0.0054 \\
Image + clinical & 0.671 & 0.869 & $+0.198$ & $-0.208$ to $+0.210$ & 0.066 & 0.33 \\
Image + clinical + radiomics & 0.619 & 0.916 & $+0.297$ & $-0.228$ to $+0.221$ & 0.0067 & 0.054 \\
Cascade & 0.742 & 0.857 & $+0.116$ & $-0.197$ to $+0.190$ & 0.40 & 0.79 \\
Multi-scale & 0.484 & 0.804 & $+0.321$ & $-0.261$ to $+0.252$ & 0.013 & 0.093 \\
Trinity & 0.563 & 0.788 & $+0.225$ & $-0.233$ to $+0.225$ & 0.053 & 0.32 \\
SegMamba-v3 & 0.755 & 0.892 & $+0.137$ & $-0.202$ to $+0.201$ & 0.26 & 0.77 \\
Leaky reference & 0.870 & 0.969 & $+0.099$ & $-0.126$ to $+0.124$ & 0.17 & 0.69 \\
Leaky reference + radiomics & 0.918 & 0.975 & $+0.057$ & $-0.101$ to $+0.103$ & 0.54 & 0.79 \\
\midrule
\multicolumn{7}{l}{\emph{Late fusion (XGBoost)}} \\
Age + sex & 0.562 & 0.794 & $+0.232$ & $-0.221$ to $+0.198$ & 0.029 & 0.17 \\
Radiomics & 0.764 & 0.771 & $+0.007$ & $-0.236$ to $+0.235$ & 0.95 & 1 \\
Image & 0.588 & 0.755 & $+0.167$ & $-0.234$ to $+0.223$ & 0.16 & 0.64 \\
Image + clinical & 0.657 & 0.832 & $+0.175$ & $-0.215$ to $+0.202$ & 0.10 & 0.51 \\
Clinical + radiomics & 0.847 & 0.835 & $-0.012$ & $-0.187$ to $+0.183$ & 0.91 & 1 \\
Image + clinical + radiomics & 0.791 & 0.846 & $+0.056$ & $-0.189$ to $+0.183$ & 0.56 & 1 \\
\bottomrule
\end{tabular}
\end{table}

Every end-to-end model scored lower on the validation split than on the test split for diagnosis
(Table~\ref{tab:valtest}), and the image-only model's validation AUC of 0.381 was below chance. The two splits have
nearly the same composition: 9 glioblastomas, 10--11 IDH-mutant astrocytomas, 3--4 IDH-wildtype astrocytomas and 2
oligodendrogliomas each, with median ages of 52 and 47 years. The checkpoint was selected on validation loss, which
should favour validation, not test. For the best end-to-end model the validation diagnosis balanced accuracy at the
selected epoch was 0.29, barely above chance, against 0.74 on test.

Chance explains much of this. Merely re-assigning patients between two such splits gives gaps with a 95\% range of
about $\pm 0.2$ for most diagnosis models and up to $\pm 0.35$ for LGG/HGG, with no change to any model. Most observed
gaps lie inside that range, and none of the 48 late-fusion gaps survived correction in any task. Gaps favoured the test
split in 24 of 27 end-to-end comparisons, 22 of 24 for XGBoost and 15 of 24 for logistic regression. These share
patients and are not independent, but together they suggest that the test split is somewhat easier.

Chance does not explain everything. The image-only diagnosis gap of $+0.418$ fell outside the chance range and stayed
significant even when corrected across all 75 gap tests (adjusted $p = 0.045$). The image + clinical + radiomics gap
was borderline (adjusted $p = 0.054$). On the same patients and inputs, the end-to-end gap was larger than the
late-fusion gap for image only ($+0.252$, adjusted $p = 0.0051$) and image + clinical + radiomics ($+0.241$, adjusted
$p = 0.037$), but not for image + clinical ($+0.023$, $p = 0.77$). Image only pointed the same way for grade ($+0.298$,
adjusted $p = 0.11$) and LGG/HGG ($+0.273$, adjusted $p = 0.098$). Part of the instability therefore belongs to the
end-to-end models rather than to the patients. Mask quality does not explain it, since the backbone's Dice was higher on
the validation patients (0.885) than on the test patients (0.791).
\todo{Re-scoring on Colab: \texttt{run("cls explain")} writes \texttt{predictions\_val.csv} from the saved checkpoints,
through the test code path, and \texttt{paper\_tables/validation\_recheck} compares it with the scores above. If they
agree, say so here: the low validation scores are real. If not, correct this table and section.}

Either way, a ranking of models by one such split is not reliable~\cite{varoquaux2018cv}.

\subsection{Error analysis}
High-grade recall was 0.7--0.8 for every model because the same patients were missed. Two glioblastomas were called
low-grade by all nine end-to-end models, including the leaky references:
\begin{itemize}
  \item In one (38 years, biopsy only, 3.7~mL tumour), the backbone's predicted mask was a 150-voxel false positive with
        Dice 0, so both the crop and the radiomics came from the wrong region.
  \item The other was segmented correctly, but its T2 appearance was that of a lower-grade tumour. Knowing it was
        IDH-wildtype, the leaky models called it an IDH-wildtype astrocytoma.
\end{itemize}
A grade-4 IDH-mutant astrocytoma (29 years) received the correct diagnosis from eight of nine models, but six of the
nine called it low grade.

Across the 25 classification test patients, the backbone missed three tumours completely (Dice 0; tumour volumes
1.6--3.7~mL) and over-segmented a fourth by a factor of 2.5. About 16\% of the classifier inputs thus rested on wrong
masks. Training patients did not show this (median Dice 0.935), because the backbone had been trained on them.
\todo{ground-truth-mask analysis (\texttt{paper\_tables/classification\_oracle}, \texttt{kfold\_oracle}): how much
of the gap these patients and the mask quality explain, end to end and in late fusion. Radiomics stability
(\texttt{radiomics\_stability}): which feature classes survive segmentation errors.}

\subsection{Explanations}
Token occlusion in the image + clinical + radiomics model showed the stage-4 image token mattering most for diagnosis
(mean drop in true-class probability 0.089, SD 0.146 across patients), ahead of clinical (0.034) and radiomics
(0.016). Removing the decoder token changed nothing (0.0001). Fusion attention agreed: of the diagnosis token's
attention, stage 4 took 0.39, radiomics 0.18, clinical 0.15 and the decoder 0.12.

Grad-CAM maps of stage 4 were coarse. In glioblastomas they lay in the heterogeneous, necrotic-looking core
(Fig.~\ref{fig:cam}), but in other patients they sat on the tumour border or outside the tumour. Decoder maps were
empty, copied the tumour mask or its complement, consistent with the classifier ignoring that token. We read these
maps as qualitative only~\cite{adebayo2018sanity} and make no claim that the classifier attends to the tumour.

\begin{figure}[t]
\centering
\includegraphics[width=0.75\textwidth]{explain/cls_image_clinical_radiomics_gradcam_panel_diagnosis.png}
\caption{Grad-CAM of the image + clinical + radiomics model for diagnosis: per class, the most confidently correct
test patient (for a class with no correct call, its most probable patient), axial slice through the tumour.
\todo{regenerate after the Colab rerun of \texttt{gradcam\_panel} (calls now flip-averaged as in the metrics);
consider dropping the decoder column.}}
\label{fig:cam}
\end{figure}

% ---------------------------------------------------------------------------------------------------------
\section{Discussion}

\paragraph{Segmentation.} The convolutional SegResNet was the best or tied model on T2, and SegMamba and SwinUNETR
were not better. This agrees with recent evidence that well-configured CNNs remain hard to beat when budgets are
matched~\cite{isensee2021nnunet,isensee2024revisited}. Self-supervised pretraining of SwinUNETR on 345 training
volumes gave no measurable gain. U-Net results depended strongly on configuration, so a weak baseline can make an
architecture look better than it is.

\paragraph{Leakage.} Near-perfect diagnosis results follow from the molecular and surgical inputs alone. This is
expected: IDH and 1p/19q define three of the four WHO~2021 classes. Inputs that are part of the label definition, or
only known after surgery, have no place in a model that claims to predict from pre-operative
MRI~\cite{kapoor2023leakage}. Where they are used, they should be reported as such.

\paragraph{What imaging adds.} Without leakage, T2 imaging carries information about diagnosis beyond clinical
variables (CV AUC 0.83 against 0.65 with XGBoost and 0.73 with logistic regression), with both classifiers and both
encodings. For grade and LGG/HGG, the gain depended on the classifier. Radiomics and deep features from the segmentation backbone were largely interchangeable.
Each added to clinical variables, deep features added nothing once radiomics were present, and radiomics were not
significantly better than deep features alone. On this cohort we thus find no reason to prefer the costlier deep
encoder, but also no evidence that it is worse. Age and sex saturated from about 32 training patients while imaging kept
gaining, so these comparisons may shift in larger cohorts. The strength of tree ensembles on engineered features on small tabular
problems~\cite{grinsztajn2022tree} fits this picture. It is not evidence that decoupled pipelines are inherently
superior: the end-to-end model with radiomics scored highest on the test split, although far lower on validation.

The subgroups add what pooled AUCs cannot show. For grade, diagnosis and LGG/HGG, radiomics and deep features stayed
above age and sex inside every age band, so the imaging signal is not merely a proxy for age. IDH status was the
exception. Age alone is a strong baseline there, because IDH-wildtype glioblastoma patients are older. Imaging gained
nothing significant over clinical variables in the pooled cohort and was level with age inside the one age band large
enough to tell, so T2 radiogenomic claims for IDH should be measured against age, also within age bands. Baselines need
care too: XGBoost on age and sex scored far below age itself, so a gain measured against a weak baseline can overstate
what imaging adds.

\paragraph{Evaluation.} With 25 test patients a balanced-accuracy CI spans about $\pm 0.2$. Re-assigning patients at
random between our two splits, of the same size and class mix, moved the diagnosis gap by about $\pm 0.2$ AUC with no
change to any model. A difference of that size between two models, or two studies, on one such split says little.
Single-split results at this size are illustrative. Model comparisons need repeated CV, several seeds, and preferably
an external cohort~\cite{varoquaux2018cv,varoquaux2022failures}. On our test split, no pre-specified difference between
the end-to-end models was significant after correction, including the radiomics token that seemed to lift diagnosis
balanced accuracy by 0.29. The largest gaps went beyond chance and were specific to end-to-end models: late fusion on
the same patients and inputs was more stable. A plausible reading is that end-to-end training on 112 patients gives
weights whose performance varies more across patients than a few selected features do, but with one seed we cannot
separate this from run-to-run variation.

The shuffled-label control shows what a common mistake costs. Selecting features before the folds gave up to 0.58 AUC
on labels with no signal, and inflated real results by up to 0.05, which is as large as the differences between feature
sets that a study like this sets out to measure. The mistake would not produce a striking result by itself, which is why
it is easy to miss without a control.

\paragraph{Recommendations.} From this study we suggest the following for small-cohort glioma work:
\begin{enumerate}
  \item Exclude any input that defines or follows from the label, and report it only as a labelled reference.
  \item Fix one patient-level split shared by all stages, with the classification test set held out of pretraining
        and segmentation training.
  \item Fit every selector and scaler inside the training folds.
  \item Always report no-imaging baselines (majority, age + sex).
  \item Report CIs and paired tests, and do not rank models whose differences are not significant.
  \item Check that each model actually uses its inputs. Silent failures, such as feature tokens that are all zeros
        or a pretrained checkpoint that fails to load, are easy to miss in multi-stage pipelines.
  \item Run a shuffled-label control through the whole pipeline. For late fusion it needs no GPU, and a leak in the
        feature steps shows up at once.
  \item Before interpreting a gap between validation and test, compare it with chance re-assignment of the patients
        between the splits. It needs only the saved predictions.
\end{enumerate}

\paragraph{Limitations.}
\begin{itemize}
  \item The data come from a single centre and use T2 only. Enhancement and necrosis are better seen on
        post-contrast T1, which likely explains part of the missed glioblastomas.
  \item The balanced cohort does not reflect clinical prevalence.
  \item End-to-end and segmentation results are from one seed, so run-to-run variation is not measured. The gap test
        holds each model's predictions fixed and says nothing about it either.
  \item The gap test and the shuffled-label control were added after the results were known, and the Holm families of
        the gap test follow this study's other comparisons rather than a prior plan. Both are exploratory. Each shuffled
        run is one 5-fold pass, noisier than the 50 folds of the real labels.
  \item The subgroup results are descriptive, one 5-fold pass per model, and several subgroups are small.
  \item The shared backbone was chosen after the segmentation benchmark and saw the classification training and
        validation images. It gives near-perfect masks on its own training patients but not on new ones, so
        classifier training crops and radiomics are cleaner for training than for test patients.
  \item One scoring rule (the derived LGG/HGG call) was changed after the first test evaluation. We report AUC as
        primary and the original values in the supplement.
  \item Grad-CAM was assessed qualitatively.
  \item There is no external validation. BraTS~\cite{menze2015brats,baid2021brats} overlaps with UCSF-PDGM; the
        Erasmus Glioma Database~\cite{vandervoort2021egd} would be a suitable external test set.
\end{itemize}

% ---------------------------------------------------------------------------------------------------------
\section{Conclusion}
In a leakage-controlled benchmark on T2-weighted MRI:
\begin{itemize}
  \item a convolutional segmenter was at least as good as state-space and transformer models, and self-supervised
        pretraining did not help;
  \item near-perfect diagnosis accuracy was reachable only with molecular inputs that define the label;
  \item the imaging signal for diagnosis was reachable through either hand-crafted radiomics or deep features, and deep
        features added nothing once radiomics were present;
  \item for grade, diagnosis and LGG versus HGG, imaging stayed above age and sex inside age bands, while for IDH status
        it added nothing significant to clinical variables and was level with age inside age bands;
  \item on shuffled labels the pipeline scored chance, and feature selection before the folds would have inflated
        results by as much as the differences between feature sets;
  \item chance re-assignment between two 25-patient splits moved macro AUC by about $\pm 0.2$, so a single such split is
        too variable to rank models.
\end{itemize}
We release the resumable pipeline so that these checks can be applied to other small-cohort studies.

\section*{Code and data availability}
Code, configurations and the experiment protocol: \url{https://github.com/kkoutsikos/glioma-mtl}. UCSF-PDGM is
publicly available~\cite{calabrese2022ucsf}. \todo{archive a tagged release (e.g.\ Zenodo DOI) and the split file.}

\section*{Acknowledgements}
\todo{supervisor, advisors, compute. Decide whether to disclose that this work builds on the first author's
postgraduate thesis~\cite{koutsikos2026thesis}; many journals ask authors to declare overlap with a thesis.}

\bibliographystyle{unsrtnat}
\bibliography{references}

% ---------------------------------------------------------------------------------------------------------
\appendix
\section{Supplementary material}

\paragraph{First evaluation of the derived LGG/HGG call.} Under the original rule (HGG only when
$P(\text{grade 4}) > 0.5$), the test LGG/HGG balanced accuracies were: image only 0.500, image + clinical 0.850,
image + clinical + radiomics 0.850, leaky reference 0.817, leaky reference + radiomics 0.850, cascade 0.767,
multi-scale 0.500, Trinity 0.500. Under the current rule they are 0.817, 0.783, 0.750, 0.833, 0.850, 0.717, 0.833 and
0.833. SegMamba-v3 has its own head (0.783) and is unaffected.

\todo{Full tables from \texttt{paper\_tables/}: all four late-fusion classifiers (CV and hold-out), macro F1 and
accuracy, the comparison tests and the ground-truth-mask analysis, ROC curves, training curves, per-patient
segmentation results.}

\end{document}
